\section{Flight dynamics: drag, Magnus lift and spin decay}
\label{sec:flight}

% Source of truth: src/flight.py (simulate, _aero, backspin_omega, cl_linear_capped, load_params),
% plus the fixed-step RK4 copies of the same force model in src/feasibility.py (integrate_batch),
% src/robustness.py (integrate3d) and src/contact.py (_acc, _rk4, free_flight).
% Parameter values: research/params_lit.json (the default for new runs; see research/LITERATURE_SYNTHESIS.md).
% Every worked number below was computed on 2026-10-30 by running the repo code with the system python
% after fl.set_aero_path(fl.LIT_PARAMS_PATH).

Once POLLEN leaves the shooter, nothing on the robot can change its path any more.
Three things act on the ball in the air: gravity, air drag and the Magnus force from spin.
This section derives, step by step, exactly what the flight solver \code{src/flight.py} computes.
The goal is to answer two practical questions.
How large is the spin force compared with gravity and drag?
And does adding more backspin keep helping?
In this model the answer to the second question is no: above a certain spin the lift coefficient
stops growing, and extra spin changes nothing.

Every aerodynamic coefficient in this section is taken from an assumed, literature-based range, not
from a POLLEN measurement.
The ranges come from \code{research/params\_lit.json}.
They are based on pickleball and smooth-sphere data, but they are not proven bounds: no one has yet
measured drag, lift or spin decay on a real POLLEN ball.

% ---------------------------------------------------------------------------------------------
\subsection{Frame, state and parameters}
\label{sec:flight-frame}

\paragraph{Motivation.}
Before writing any force we need a coordinate frame and a list of what the solver tracks.
Sign errors in the Magnus force are the most common bug in spin models, so the frame must be fixed first.

\paragraph{Frame.}
The world frame is right-handed.
The axis $x$ is horizontal and points toward the target, $y$ is horizontal and sideways, and $z$ points up.
The unit vectors are $\hat{\vect{x}}$, $\hat{\vect{y}}$ and $\hat{\vect{z}}$.
All quantities are in SI units (m, s, kg, rad).

\paragraph{State.}
The solver integrates a nine-component state: the position of the ball centre $\vect{r}$ (\unit{m}),
its velocity $\vect{v}$ (\unit{m/s}) and its spin (angular velocity) vector $\vect{\omega}$ (\unit{rad/s}).
The ball is treated as a rigid sphere of mass $m$ (\unit{kg}) and outer radius $R$ (\unit{m}).

\paragraph{Parameters.}
Table~\ref{tab:flight-params} lists the physical constants read by \code{load\_params}.
The acceleration of gravity is $g$ (\unit{m/s\textsuperscript{2}}), the air density is $\rho$ (\unit{kg/m\textsuperscript{3}}) and the
kinematic viscosity of air is $\nu$ (\unit{m\textsuperscript{2}/s}).
The frontal (cross-sectional) area of the ball is
\begin{equation}
  A = \pi R^2 ,
  \label{eq:flight-area}
\end{equation}
in \unit{m\textsuperscript{2}}.
With the nominal radius, $A = 3.959\times10^{-3}$\,\unit{m\textsuperscript{2}}\derived.

\begin{table}[htbp]
  \centering
  \caption{Physical constants used by the flight solver (\code{research/params\_lit.json}).
    Only the nominal values of $\rho$, $\nu$, $m$ and $R$ are read by \code{load\_params}; the bounds
    are used by the robustness sweeps.}
  \label{tab:flight-params}
  \small
  \begin{tabular}{@{}lllll@{}}
    \toprule
    Symbol & Meaning & Unit & Nominal & Bounds \\
    \midrule
    $g$    & standard gravity & \unit{m/s\textsuperscript{2}} & 9.80665\confirmed & -- \\
    $\rho$ & air density & \unit{kg/m\textsuperscript{3}} & 1.2\assumed & 1.0\assumed\ to 1.25\assumed \\
    $\nu$  & kinematic viscosity & \unit{m\textsuperscript{2}/s} & $1.5\times10^{-5}$\assumed & $1.45\times10^{-5}$\assumed\ to $1.6\times10^{-5}$\assumed \\
    $m$    & ball mass & \unit{kg} & 0.02495\assumed & 0.0235\assumed\ to 0.0275\assumed \\
    $R$    & outer radius & \unit{m} & 0.0355\confirmed & 0.035\assumed\ to 0.036\assumed \\
    \bottomrule
  \end{tabular}
\end{table}

% Provenance note: m = 24.95 g is the AndyMark listing (0.055 lb), grade C vendor data in params_lit.json,
% not the official CAD or game manual, so it is tagged \assumed (as in the contact section). The radius
% 0.0355 m is the manual's "approximately 2.8 in (7.1 cm)", hence \confirmed.
% Code note: src/flight.py sets AERO_PARAMS_PATH to data/aero_params.json unless the environment variable
% PM001_AERO_PATH is set. Only the command-line entry points of feasibility.py, robustness.py and contact.py
% call set_aero_path(LIT_PARAMS_PATH). A bare fl.load_params() or fl.default_params() (and the unit tests
% test_hive / test_contact) therefore still read the older first-pass file, although DECISIONS.md makes
% params_lit.json the default for new runs. All numbers in this section use params_lit.json explicitly.

\paragraph{Assumptions.}
\begin{itemize}
  \item The world frame is inertial (no Coriolis or centrifugal terms) and the floor is flat.
  \item The ball is a rigid sphere of fixed $m$ and $R$; it does not deform in flight.
  \item $\rho$, $\nu$ and $g$ are constant during a shot.
\end{itemize}

% ---------------------------------------------------------------------------------------------
\subsection{Equations of motion}
\label{sec:flight-eom}

\paragraph{Motivation.}
Newton's second law gives the translation; a simple decay law gives the spin.
The only subtlety is wind: aerodynamic forces depend on how the ball moves \emph{through the air},
not over the floor.

\paragraph{Air-relative velocity.}
Let $\vect{u}$ (\unit{m/s}) be the wind velocity.
The air-relative velocity is
\begin{equation}
  \vect{v}_{\mathrm{rel}} = \vect{v} - \vect{u} .
  \label{eq:flight-vrel}
\end{equation}
Both aerodynamic forces below use $\vect{v}_{\mathrm{rel}}$; gravity does not.
The code accepts a constant $\vect{u}$ or a function $\vect{u}(t,\vect{r})$.
The default is still air, $\vect{u}=\vect{0}$, so then $\vect{v}_{\mathrm{rel}}=\vect{v}$.

\paragraph{Newton's law and the spin law.}
Gravity is $m\vect{g}$ with $\vect{g} = -g\,\hat{\vect{z}}$.
Adding the drag force $\vect{F}_D$ and the Magnus force $\vect{F}_M$ (both in \unit{N}) gives
\begin{align}
  \frac{d\vect{r}}{dt} &= \vect{v}, \label{eq:flight-eom-r}\\
  m\,\frac{d\vect{v}}{dt} &= -m g\,\hat{\vect{z}} + \vect{F}_D(\vect{v}_{\mathrm{rel}}) + \vect{F}_M(\vect{v}_{\mathrm{rel}},\vect{\omega}), \label{eq:flight-eom-v}\\
  \frac{d\vect{\omega}}{dt} &= -k\,\vect{\omega}. \label{eq:flight-eom-w}
\end{align}
Here $k$ (\unit{1/s}) is the spin-decay rate (Section~\ref{sec:flight-decay}).
In \code{\_make\_rhs} the right-hand side is exactly this:
the acceleration is $(\vect{F}_D+\vect{F}_M)/m$ with $g$ subtracted from the $z$ component,
and each spin component is multiplied by $-k$.
No other torque acts on the ball, so the spin axis never turns.

\paragraph{Assumptions.}
\begin{itemize}
  \item The ball is a point mass for translation; only its spin vector is tracked for rotation.
  \item Buoyancy and the added mass of the surrounding air are neglected.
    The air displaced by the outer envelope has a mass of $2.25\times10^{-4}$\,\unit{kg}\derived,
    which is $0.9$\,\%\derived\ of $m$.
  \item The wind is prescribed; the ball's own wake and the robot's airflow are ignored.
\end{itemize}

% ---------------------------------------------------------------------------------------------
\subsection{Drag}
\label{sec:flight-drag}

\paragraph{Motivation.}
Drag is the largest aerodynamic force on POLLEN for almost every parameter set in this model
(the one exception, high lift with low drag, is in Section~\ref{sec:flight-scale}).
It slows the ball and shortens the shot.

\paragraph{Derivation.}
The dynamic pressure $q$ (\unit{Pa}) is the kinetic energy of the oncoming air per unit volume.
The drag magnitude is the dynamic pressure times the frontal area times a dimensionless drag coefficient
$C_d$.
Drag points straight against the air-relative velocity, so we multiply by the unit vector
$-\vect{v}_{\mathrm{rel}}/|\vect{v}_{\mathrm{rel}}|$:
\begin{align}
  q &= \tfrac12\,\rho\,|\vect{v}_{\mathrm{rel}}|^2, \label{eq:flight-q}\\
  |\vect{F}_D| &= q\,C_d\,A, \label{eq:flight-FDmag}\\
  \vect{F}_D &= -|\vect{F}_D|\,\frac{\vect{v}_{\mathrm{rel}}}{|\vect{v}_{\mathrm{rel}}|}
             = -\tfrac12\,\rho\,C_d\,A\,|\vect{v}_{\mathrm{rel}}|\,\vect{v}_{\mathrm{rel}}. \label{eq:flight-drag}
\end{align}
In the last step one factor of $|\vect{v}_{\mathrm{rel}}|$ cancels.
The form $|\vect{v}_{\mathrm{rel}}|\,\vect{v}_{\mathrm{rel}}$ is what \code{\_aero} computes
(\code{kd = -q*cd*vmag}, then \code{kd*vx}, \dots; note that the code's variable \code{q} is
$\tfrac12\rho A$, not the dynamic pressure $q$ of Eq.~\eqref{eq:flight-q}).
It has no division, so drag goes smoothly to zero as $|\vect{v}_{\mathrm{rel}}|\to 0$.

\paragraph{Values.}
$C_d$ is constant during a run.
The bounds are $C_d = 0.33$\assumed\ (low), $0.5$\assumed\ (nominal) and $0.9$\assumed\ (high).
The low end comes from outdoor pickleball free flight; the high end covers a spinning perforated ball,
because smooth-sphere data show that spin raises drag.
The solver also accepts a function $C_d(\mathrm{Re})$ or $C_d(\mathrm{Re},S)$, but no sweep uses one.
% Code vs params_lit: params_lit.json lists an optional spin term Cd = Cd0 + dCd*min(S/0.6, 1). load_params
% returns only the constant low/nominal/high value, so that term is not active in any run.

\paragraph{Worked example.}
Take a launch at $|\vect{v}| = 7$\,\unit{m/s}\assumed\ and elevation $60^\circ$\assumed\ in still air,
so $\vect{v} = (3.5\assumed, 0, 6.062\derived)$\,\unit{m/s}.
With the nominal values, $\tfrac12\rho A = 2.3755\times10^{-3}$\,\unit{kg/m}\derived, so
$|\vect{F}_D| = 2.3755\times10^{-3}\derived \times 0.5\assumed \times 49\derived = 0.0582$\,\unit{N}\derived.
\code{aero\_forces} returns $\vect{F}_D = (-0.02910, 0, -0.05040)$\,\unit{N}\predicted, which is the
same magnitude pointing exactly against $\vect{v}$.
The weight is $mg = 0.2447$\,\unit{N}\derived, so drag is $0.238$\derived\ of the weight at this speed.

\paragraph{Assumptions.}
\begin{itemize}
  \item $C_d$ is constant along the path; no Reynolds-number or spin dependence is active.
  \item No ``knuckle'' side force from the holes is modelled. \code{params\_lit.json} lists a
    non-Magnus hole force as a scatter source, but the solver does not include it.
\end{itemize}

% ---------------------------------------------------------------------------------------------
\subsection{Reynolds number and spin parameter}
\label{sec:flight-Re-S}

\paragraph{Motivation.}
Two dimensionless numbers describe the flow around a spinning ball.
The Reynolds number says which flow regime we are in; the spin parameter says how fast the surface
moves compared with the air.

\paragraph{Reynolds number.}
With the ball diameter $D = 2R$ (\unit{m}),
\begin{equation}
  \mathrm{Re} = \frac{|\vect{v}_{\mathrm{rel}}|\,2R}{\nu}.
  \label{eq:flight-Re}
\end{equation}
For speeds of $2$\assumed, $5$\assumed, $7$\assumed\ and $10$\,\unit{m/s}\assumed\ this gives
$\mathrm{Re} = 9.47\times10^{3}$\derived, $2.37\times10^{4}$\derived, $3.31\times10^{4}$\derived\
and $4.73\times10^{4}$\derived.
This is below the smooth-sphere drag crisis, but it overlaps the band where smooth-sphere lift
data show a dip (the ``lift crisis'' recorded in \code{params\_lit.json}).
In the default model $\mathrm{Re}$ is only a diagnostic: $C_d$ is constant and $C_L$ depends on $S$ alone.
It is reported so that a run can be checked against the range of the literature data.

\paragraph{Spin parameter.}
A point on the equator of the ball moves at speed $|\vect{\omega}|R$ relative to the centre.
The spin parameter is the ratio of that surface speed to the airspeed:
\begin{equation}
  S = \frac{|\vect{\omega}|\,R}{|\vect{v}_{\mathrm{rel}}|}.
  \label{eq:flight-S}
\end{equation}
$S$ has no units.
$S = 1$ means the surface speed equals the airspeed, which is what an ideal no-slip hood gives at the exit.
Note that the code uses the full spin magnitude $|\vect{\omega}|$, not only the part of $\vect{\omega}$
perpendicular to $\vect{v}_{\mathrm{rel}}$.

\paragraph{Regularisation.}
Equation~\eqref{eq:flight-S} divides by zero when the ball is at rest relative to the air.
The code therefore uses
\begin{equation}
  S = \frac{|\vect{\omega}|\,R}{\sqrt{|\vect{v}_{\mathrm{rel}}|^2 + \varepsilon_v^2}},
  \label{eq:flight-S-reg}
\end{equation}
with the regularisation speed $\varepsilon_v = 10^{-6}$\,\unit{m/s}\assumed.
For any airspeed above about $10^{-3}$\,\unit{m/s}\assumed\ the difference is far below measurement accuracy.
% HANDOFF 4A writes S = |omega| R / max(|v_rel|, eps). The code uses sqrt(|v_rel|^2 + eps^2), as here.

\paragraph{Worked example.}
For the launch above, an exit spin parameter $S_0 = 1.06$\assumed\ (the nominal hood value in
\code{params\_lit.json}) needs
$|\vect{\omega}_0| = S_0|\vect{v}|/R = 1.06\assumed\times 7\assumed/0.0355\confirmed = 209.0$\,\unit{rad/s}\derived,
which is $1996$\,\unit{rpm}\derived.
The Reynolds number is $3.31\times10^{4}$\derived.

% ---------------------------------------------------------------------------------------------
\subsection{The Magnus force}
\label{sec:flight-magnus}

\paragraph{Motivation.}
A spinning ball drags the air near its surface around with it.
The flow is deflected to one side, and the ball is pushed to the other side.
This sideways push is the Magnus force.
It is the only way spin changes the flight path in this model.

\paragraph{Magnitude and direction.}
The Magnus force has the same dynamic-pressure form as drag, with a lift coefficient $C_L$ (no units)
in place of $C_d$.
Its direction is perpendicular to both the spin axis and the airflow, given by the cross product
$\vect{\omega}\times\vect{v}_{\mathrm{rel}}$:
\begin{align}
  |\vect{F}_M| &= \tfrac12\,\rho\,C_L\,A\,|\vect{v}_{\mathrm{rel}}|^2, \label{eq:flight-FMmag}\\
  \hat{\vect{n}}_M &= \frac{\vect{\omega}\times\vect{v}_{\mathrm{rel}}}{|\vect{\omega}\times\vect{v}_{\mathrm{rel}}|}, \label{eq:flight-nM}\\
  \vect{F}_M &= |\vect{F}_M|\,\hat{\vect{n}}_M
             = \tfrac12\,\rho\,C_L(S)\,A\,|\vect{v}_{\mathrm{rel}}|^2\,
               \frac{\vect{\omega}\times\vect{v}_{\mathrm{rel}}}{|\vect{\omega}\times\vect{v}_{\mathrm{rel}}|}. \label{eq:flight-FM}
\end{align}
The unit vector $\hat{\vect{n}}_M$ carries only the direction, so the spin rate enters the force
only through $C_L(S)$.
This fact matters a great deal in Section~\ref{sec:flight-lift}.

\paragraph{Regularised form used in the code.}
To keep the direction finite when $\vect{\omega}\times\vect{v}_{\mathrm{rel}}\to\vect{0}$, \code{\_aero}
computes
\begin{equation}
  \vect{F}_M = \tfrac12\,\rho\,C_L(S)\,A\,|\vect{v}_{\mathrm{rel}}|^2\,
  \frac{\vect{\omega}\times\vect{v}_{\mathrm{rel}}}
       {\sqrt{|\vect{\omega}\times\vect{v}_{\mathrm{rel}}|^2 + \varepsilon_v^2|\vect{\omega}|^2}} ,
  \label{eq:flight-FM-reg}
\end{equation}
and sets $\vect{F}_M = \vect{0}$ when $\vect{\omega}=\vect{0}$, when $\vect{v}_{\mathrm{rel}}=\vect{0}$
or when no lift law is given.
Because of the factor $|\vect{v}_{\mathrm{rel}}|^2$ and the bounded $C_L$, the force vanishes smoothly
at zero airspeed; there is no fictitious force from the $1/|\vect{v}_{\mathrm{rel}}|$ in $S$.

\paragraph{Magnus does no work in still air.}
A useful property follows from the cross product.
For any vectors, $(\vect{\omega}\times\vect{v}_{\mathrm{rel}})\cdot\vect{v}_{\mathrm{rel}} = 0$, so
\begin{equation}
  \vect{F}_M\cdot\vect{v}_{\mathrm{rel}} = 0 .
  \label{eq:flight-FM-perp}
\end{equation}
In still air $\vect{v}_{\mathrm{rel}}=\vect{v}$, so the power $\vect{F}_M\cdot\vect{v}$ is zero.
The Magnus force bends the path but cannot speed the ball up or slow it down.
Only gravity and drag change the speed.

\paragraph{Worked example.}
For the launch above ($S_0 = 1.06$\assumed, nominal set, so $C_L = 0.25$\derived; see below),
\code{aero\_forces} returns $\vect{F}_M = (-0.02520, 0, 0.01455)$\,\unit{N}\predicted\ with
$|\vect{F}_M| = 0.0291$\,\unit{N}\predicted.
Check: $2.3755\times10^{-3}\derived\times0.25\derived\times49\derived = 0.0291$\,\unit{N}\derived.
The dot product with $\vect{v} = (3.5, 0, 6.062)$ is
$-0.02520\times3.5 + 0.01455\times6.062 = 0.0000$\derived\ (to four decimals), as
Eq.~\eqref{eq:flight-FM-perp} requires.

\paragraph{Assumptions.}
\begin{itemize}
  \item $C_L$ depends on $S$ only, and $S$ uses $|\vect{\omega}|$, not its component perpendicular to
    $\vect{v}_{\mathrm{rel}}$. For pure backspin in a vertical plane the two are equal.
  \item The direction rule $\hat{\vect{n}}_M = \operatorname{unit}(\vect{\omega}\times\vect{v}_{\mathrm{rel}})$
    holds for every $S$: the model has no reverse (negative) Magnus effect.
    The literature judges reverse Magnus unlikely at POLLEN's Reynolds numbers (the documented
    negative-lift windows are at $\mathrm{Re} \ge 4.6\times10^{4}$\assumed\ and $S < 0.6$\assumed),
    but does not exclude it, because perforated-ball lift data stop at $S \approx 0.56$\assumed\ in
    backspin (Van Wyk et al.\ 2026) and $S \approx 0.65$\assumed\ in topspin (Lindsey 2025, an unreviewed
    web article).
    % Fixed: the draft said "unlikely at S >= 0.5"; params_lit.json Cl.negative_magnus states the Re and S
    % windows above instead.
\end{itemize}

% ---------------------------------------------------------------------------------------------
\subsection{The lift law and its saturation}
\label{sec:flight-lift}

\paragraph{Motivation.}
How $C_L$ grows with spin decides whether more spin helps.
For POLLEN the shape of $C_L(S)$ is the largest unknown in the whole project.

\paragraph{Linear-capped law.}
\code{cl\_linear\_capped} implements
\begin{equation}
  C_L(S) = \min\bigl(c_1 S,\; C_{L,\max}\bigr), \qquad S \ge 0,
  \label{eq:flight-cl}
\end{equation}
where $c_1$ (no units) is the initial slope and $C_{L,\max}$ (no units) is the ceiling.
The law rises linearly from zero and then stays flat.
It has no offset, so $C_L = 0$ when $S = 0$.

\paragraph{Saturation spin.}
The two branches meet where $c_1 S = C_{L,\max}$.
Solving for $S$ gives the saturation spin parameter
\begin{equation}
  S_{\mathrm{sat}} = \frac{C_{L,\max}}{c_1}.
  \label{eq:flight-ssat}
\end{equation}
For $S \ge S_{\mathrm{sat}}$ the code returns $C_{L,\max}$ exactly
(\code{return v if v < cl\_max else cl\_max}).
Table~\ref{tab:flight-clsets} gives the four sets in \code{params\_lit.json}.

\begin{table}[htbp]
  \centering
  \caption{Assumed lift sets in \code{research/params\_lit.json}, based on pickleball data at
    $S \le 0.56$\assumed\ and on smooth-sphere data. No perforated-ball lift has been measured at the
    $S \approx 1$\assumed\ to $2.5$\predicted\ of a POLLEN shot, so the sets are not proven bounds and
    not POLLEN measurements. $S_{\mathrm{sat}}$ is from Eq.~\eqref{eq:flight-ssat}.}
  \label{tab:flight-clsets}
  \begin{tabular}{@{}lllll@{}}
    \toprule
    Set & $c_1$ & $C_{L,\max}$ & $S_{\mathrm{sat}}$ & Basis (summary) \\
    \midrule
    floor   & 0.02\assumed & 0.04\assumed & 2.00\derived & pickleball fit, Magnus-only reading \\
    low     & 0.2\assumed  & 0.15\assumed & 0.75\derived & pickleball slope; smooth-sphere lift dip \\
    nominal & 0.35\assumed & 0.25\assumed & 0.714\derived & between pickleball and smooth sphere \\
    high    & 0.8\assumed  & 0.4\assumed  & 0.50\derived & smooth-sphere slope and plateau \\
    \bottomrule
  \end{tabular}
\end{table}
% Code vs docs: docs/02_aero.md tabulates the older data/aero_params.json low set (c1 0.15, cap 0.10,
% S_sat 0.67) and has no floor set. params_lit.json (followed here) has low = 0.2/0.15 and floor = 0.02/0.04.
% docs/02_aero.md section 5 also says "above S ~ 0.5-0.7, extra spin buys nothing under any set"; with
% params_lit.json that is false for the floor set (S_sat = 2.0) and the low set saturates at 0.75.

\paragraph{How $S$ changes along the path.}
With constant spin decay (Eq.~\eqref{eq:flight-decay-sol} below) the spin parameter along the path is
\begin{equation}
  S(t) = \frac{|\vect{\omega}_0|\,e^{-kt}\,R}{|\vect{v}_{\mathrm{rel}}(t)|}
       = S_0\,e^{-kt}\,\frac{|\vect{v}_{\mathrm{rel}}(0)|}{|\vect{v}_{\mathrm{rel}}(t)|}.
  \label{eq:flight-S-t}
\end{equation}
The ball slows on the way up, so $S$ \emph{rises} toward the apex even though the spin itself decays.
In the worked shot below, $S$ grows from $1.06$\assumed\ at launch to $2.53$\predicted\ at the apex and
falls back to $1.08$\predicted\ at the floor.

\paragraph{Why spin above $S_{\mathrm{sat}}$ changes nothing in this model.}
Look at where $\vect{\omega}$ enters the equations of motion.
The translational equation~\eqref{eq:flight-eom-v} depends on $\vect{\omega}$ only through
(i) $C_L(S)$ and (ii) the unit vector $\hat{\vect{n}}_M$.
The spin law~\eqref{eq:flight-eom-w} scales $\vect{\omega}$ but never turns it, so $\hat{\vect{n}}_M$
depends on the spin \emph{direction}, which is fixed at launch, and not on the spin rate.
Now suppose $S(t) \ge S_{\mathrm{sat}}$ for the whole flight.
Then $C_L = C_{L,\max}$ at every instant, and the Magnus force becomes
\begin{equation}
  \vect{F}_M = \tfrac12\,\rho\,C_{L,\max}\,A\,|\vect{v}_{\mathrm{rel}}|^2\,\hat{\vect{n}}_M
  \qquad (S \ge S_{\mathrm{sat}}),
  \label{eq:flight-FM-sat}
\end{equation}
which contains no spin rate at all.
Two launches with the same $\vect{v}_0$ and spin axis but different spin rates, both above
saturation for the whole flight, therefore obey identical equations from identical initial
positions and velocities.
They follow the same path.
For the same reason the decay rate $k$ has no effect on the path while $S$ stays above $S_{\mathrm{sat}}$.

\textbf{This is a model assumption, not a measured property of POLLEN.}
The cap was chosen because smooth-sphere and baseball data flatten out and because no data exist for a
perforated ball at $S \approx 1$ to $2$.
A real ball might keep gaining lift, might show a dip (the \code{dip\_variant} table in
\code{params\_lit.json}, which \code{src/flight.py} does not read), and might gain drag with spin,
which the constant-$C_d$ model leaves out.
So the model cannot support the claim that more spin always helps, and it does not make that claim.

\paragraph{Worked example.}
Release from $(0, 0, 0.4)$\,\unit{m}\assumed\ at $7$\,\unit{m/s}\assumed\ and $60^\circ$\assumed,
nominal set, $k = 0.03$\,\unit{1/s}\assumed.
Running \code{simulate} with $S_0 = 1.0$\assumed, $1.06$\assumed, $1.5$\assumed\ and $2.0$\assumed\
gives the same apex at $x = 1.846$\,\unit{m}\predicted, $z = 2.155$\,\unit{m}\predicted\ and the same
landing point $x = 3.764$\,\unit{m}\predicted.
The largest position difference between the $S_0 = 1.0$ and $S_0 = 2.0$ paths is
$0.0$\,\unit{m}\predicted\ in double precision, because the smallest $S$ on the $S_0 = 1.0$ path is
$1.0$\predicted, already above $S_{\mathrm{sat}} = 0.714$\derived.
Changing $k$ to $0.005$\assumed\ or $0.15$\,\unit{1/s}\assumed\ also leaves the path unchanged.
Table~\ref{tab:flight-setfx} compares the sets against the same launch without spin.
Only the floor set responds to the extra spin: its $S_{\mathrm{sat}} = 2.0$\derived\ lies above
$S_0 = 1.06$, so with $S_0 = 1.06$ it is below its cap for most of the path (it saturates only around the
apex, from $t = 0.39$\predicted\ to $0.79$\,\unit{s}\predicted\ of a $1.25$\,\unit{s}\predicted\ flight,
where $S$ peaks at $2.38$\predicted), while with $S_0 = 2.0$ it is capped for the whole flight.

\begin{table}[htbp]
  \centering
  \caption{Change caused by backspin relative to the same launch with no spin
    ($7$\,\unit{m/s}\assumed, $60^\circ$\assumed, release height $0.4$\,\unit{m}\assumed, $C_d = 0.5$\assumed,
    $k = 0.03$\,\unit{1/s}\assumed). $\Delta z_{\mathrm{apex}}$ is the apex height gain and
    $\Delta x_{\mathrm{land}}$ the landing-point shift, in cm. One model run per cell; these are model
    sensitivities, not repeatable test results.}
  \label{tab:flight-setfx}
  \begin{tabular}{@{}lllll@{}}
    \toprule
    & \multicolumn{2}{c}{$\Delta z_{\mathrm{apex}}$ (cm)} & \multicolumn{2}{c}{$\Delta x_{\mathrm{land}}$ (cm)} \\
    \cmidrule(lr){2-3}\cmidrule(l){4-5}
    Set & $S_0{=}1.06$ & $S_0{=}2.0$ & $S_0{=}1.06$ & $S_0{=}2.0$ \\
    \midrule
    floor   & $+0.74$\predicted & $+1.06$\predicted & $-0.18$\predicted & $-0.67$\predicted \\
    low     & $+3.98$\predicted & $+3.98$\predicted & $-2.86$\predicted & $-2.86$\predicted \\
    nominal & $+6.60$\predicted & $+6.60$\predicted & $-5.29$\predicted & $-5.29$\predicted \\
    high    & $+10.47$\predicted & $+10.47$\predicted & $-9.72$\predicted & $-9.72$\predicted \\
    \bottomrule
  \end{tabular}
\end{table}

The landing shift is negative: on this steep shot backspin makes the ball land slightly \emph{shorter}.
Section~\ref{sec:flight-axis} shows why: while the ball rises, the Magnus force points up \emph{and back}.

\paragraph{Assumptions.}
\begin{itemize}
  \item $C_L$ follows Eq.~\eqref{eq:flight-cl} with one of the four sets; no Reynolds-number dependence,
    no lift dip, no offset.
  \item The cap $C_{L,\max}$ holds for all $S$ up to the largest value reached in flight
    (\code{params\_lit.json} gives a valid range of $0$ to $2$\assumed; apex values above that, as in the
    example, are an extrapolation of the flat cap).
\end{itemize}

% ---------------------------------------------------------------------------------------------
\subsection{Spin decay}
\label{sec:flight-decay}

\paragraph{Motivation.}
Air friction on the surface slows the spin during flight.
The model uses the simplest law that captures this.

\paragraph{Derivation.}
Equation~\eqref{eq:flight-eom-w} acts on each component separately: $d\omega_i/dt = -k\,\omega_i$.
Separating variables, $d\omega_i/\omega_i = -k\,dt$, and integrating from $0$ to $t$ gives
$\ln(\omega_i(t)/\omega_{i,0}) = -kt$, so
\begin{equation}
  \vect{\omega}(t) = \vect{\omega}_0\,e^{-kt}.
  \label{eq:flight-decay-sol}
\end{equation}
All components shrink by the same factor, so the spin axis stays fixed in space.
The time for the spin to halve is
\begin{equation}
  t_{1/2} = \frac{\ln 2}{k}.
  \label{eq:flight-halflife}
\end{equation}
\code{flight.simulate} and \code{contact.free\_flight} integrate Eq.~\eqref{eq:flight-eom-w} as part of the
state; \code{feasibility.integrate\_batch} and \code{robustness.integrate3d} use the closed form
\eqref{eq:flight-decay-sol} directly.

\paragraph{Values.}
\begin{table}[htbp]
  \centering
  \caption{Spin-decay bounds from \code{params\_lit.json} and their effect over a
    $1.2$\,\unit{s}\assumed\ flight.}
  \label{tab:flight-decay}
  \begin{tabular}{@{}llll@{}}
    \toprule
    Bound & $k$ (\unit{1/s}) & Spin lost in $1.2$\,\unit{s}\assumed & $t_{1/2}$ (\unit{s}) \\
    \midrule
    low     & 0.005\assumed & 0.6\,\%\derived  & 139\derived \\
    nominal & 0.03\assumed  & 3.5\,\%\derived  & 23.1\derived \\
    high    & 0.15\assumed  & 16.5\,\%\derived & 4.62\derived \\
    \bottomrule
  \end{tabular}
\end{table}
% Code vs docs: docs/02_aero.md section 6 quotes the older aero_params.json range k = 0.01/0.05/0.30 1/s
% and losses of 1/6/30 % over 1.2 s. params_lit.json (followed here) uses 0.005/0.03/0.15 1/s.

\paragraph{Where $k$ comes from (literature model, not in the code).}
\code{params\_lit.json} records a torque model in which the spin-down torque grows with speed.
With a moment coefficient $C_M = \beta S$ ($\beta$ has no units), the torque is
$\tau = -R\,\rho\,A\,C_M\,|\vect{v}|^2$ (\unit{N\,m}).
The ball's moment of inertia is $I_b = \alpha m R^2$ (\unit{kg\,m\textsuperscript{2}}).
The shape factor $\alpha$ has no units.
Substituting $S = |\vect{\omega}|R/|\vect{v}|$ and $A = \pi R^2$ into $I_b\,d|\vect{\omega}|/dt = \tau$:
\begin{align}
  \alpha m R^2\,\frac{d|\vect{\omega}|}{dt}
    &= -R\,\rho\,\pi R^2\,\beta\,\frac{|\vect{\omega}|R}{|\vect{v}|}\,|\vect{v}|^2
     = -\pi\rho\beta R^4\,|\vect{v}|\,|\vect{\omega}|, \nonumber\\
  k_{\mathrm{torque}} &= \frac{\pi\,\rho\,\beta\,R^2\,|\vect{v}|}{\alpha\,m}. \label{eq:flight-ktorque}
\end{align}
With $\rho = 1.2$\assumed, $R = 0.0355$\confirmed, $\alpha = 0.63$\derived\ and $m = 0.02495$\assumed\
this is $k_{\mathrm{torque}} = 0.302\,\beta\,|\vect{v}|$\derived.
The nominal $\beta = 0.02$\assumed\ at $5$\,\unit{m/s}\assumed\ gives $0.030$\,\unit{1/s}\derived, which is
where the nominal constant $k$ comes from.
The solver uses a \emph{constant} $k$; the speed-dependent form is not implemented.
% Code vs params_lit: params_lit.json calls the speed-dependent k "preferred"; src/flight.py only
% supports a constant spin_decay_k.

\paragraph{Assumptions.}
\begin{itemize}
  \item $k$ is constant during a flight, equal to its value near $5$\,\unit{m/s}\assumed.
  \item No torque turns the spin axis (no gyroscopic or wobble effects).
\end{itemize}

% ---------------------------------------------------------------------------------------------
\subsection{The backspin axis convention}
\label{sec:flight-axis}

\paragraph{Motivation.}
``Backspin'' must be turned into a definite vector before the solver can use it.
\code{backspin\_omega} does this, and a unit test checks that backspin on a ball moving forward
gives an upward Magnus force.

\paragraph{Definition.}
For a launch velocity $\vect{v}_0$ and a spin rate $\omega_0$ (\unit{rad/s}; positive for backspin,
negative for topspin),
\begin{equation}
  \vect{\omega}_0 = \omega_0\,\frac{\vect{v}_0\times\hat{\vect{z}}}{|\vect{v}_0\times\hat{\vect{z}}|}.
  \label{eq:flight-axis}
\end{equation}
The axis is horizontal and perpendicular to the vertical plane of motion.
It is undefined for a vertical launch; the code raises an error then.

\paragraph{Derivation for a shot along $+x$.}
Write $\vect{v}_0 = (v_0\cos\theta,\,0,\,v_0\sin\theta)$ with elevation $\theta$ (\unit{rad}).
Using $\vect{a}\times\hat{\vect{z}} = (a_y,\,-a_x,\,0)$,
\begin{equation}
  \vect{v}_0\times\hat{\vect{z}} = (0,\,-v_0\cos\theta,\,0)
  \;\Longrightarrow\;
  \vect{\omega}_0 = -\,\omega_0\,\hat{\vect{y}} \quad (\cos\theta > 0).
  \label{eq:flight-axis-x}
\end{equation}
So backspin on a ball moving toward $+x$ points along $-y$.
Now take the cross product with the velocity (still air, so $\vect{v}_{\mathrm{rel}}=\vect{v}$).
Using $\hat{\vect{y}}\times\hat{\vect{x}} = -\hat{\vect{z}}$ and $\hat{\vect{y}}\times\hat{\vect{z}} = \hat{\vect{x}}$,
\begin{align}
  \vect{\omega}_0\times\vect{v}_0
    &= -\omega_0\,\bigl(v_0\cos\theta\;\hat{\vect{y}}\times\hat{\vect{x}} + v_0\sin\theta\;\hat{\vect{y}}\times\hat{\vect{z}}\bigr)
     = \omega_0 v_0\,(-\sin\theta,\,0,\,\cos\theta), \nonumber\\
  \hat{\vect{n}}_M &= (-\sin\theta,\,0,\,\cos\theta) \qquad (\omega_0 > 0). \label{eq:flight-nM-back}
\end{align}
This is the flight direction $(\cos\theta, 0, \sin\theta)$ turned by $+90^\circ$ in the $x$--$z$ plane.
For a level shot ($\theta = 0$) it points straight up, which is the unit-tested convention.
For a rising ball ($\theta > 0$) it points up \emph{and backward}; for a falling ball it points up and
forward.
That is why backspin shortened the steep shot in Table~\ref{tab:flight-setfx}.

\paragraph{Worked example.}
For $\theta = 60^\circ$\assumed\ and $\omega_0 = 209.0$\,\unit{rad/s}\derived, the code gives
$\vect{\omega}_0 = (0, -209.0, 0)$\,\unit{rad/s}\derived\ and
$\hat{\vect{n}}_M = (-0.866, 0, 0.500)$\derived.
This matches $(-\sin 60^\circ, 0, \cos 60^\circ)$.

\paragraph{Assumptions.}
\begin{itemize}
  \item The axis is set once, from $\vect{v}_0$, and stays fixed.
    If the horizontal velocity ever reversed, the same $\vect{\omega}$ would act as topspin; the code
    flags such paths with \code{horizontal\_reversal}.
  \item No sidespin or axis tilt is used in the 2D baseline.
\end{itemize}

% ---------------------------------------------------------------------------------------------
\subsection{The 2D reduction}
\label{sec:flight-2d}

\paragraph{Motivation.}
Most sweeps run in the vertical $x$--$z$ plane, because 2D runs are far cheaper than 3D runs.
We now show that pure backspin with no sideways wind keeps the ball in that plane, and derive the
scalar equations the fast sweeps use.

\paragraph{Derivation.}
Take $y_0 = 0$, $v_{y,0} = 0$, $u_y = 0$ and $\vect{\omega} = \omega_y\,\hat{\vect{y}}$
(for backspin toward $+x$, $\omega_y < 0$).
Write $\vect{v}_{\mathrm{rel}} = (v_x, 0, v_z)$ for its components.
Then
\begin{equation}
  \vect{\omega}\times\vect{v}_{\mathrm{rel}}
  = (\omega_y v_z - 0,\; 0 - 0,\; 0 - \omega_y v_x)
  = \omega_y\,(v_z,\,0,\,-v_x),
  \label{eq:flight-2d-cross}
\end{equation}
so $|\vect{\omega}\times\vect{v}_{\mathrm{rel}}| = |\omega_y|\,|\vect{v}_{\mathrm{rel}}|$ and
\begin{equation}
  \vect{F}_M = \tfrac12\,\rho\,C_L\,A\,|\vect{v}_{\mathrm{rel}}|\,\operatorname{sgn}(\omega_y)\,(v_z,\,0,\,-v_x).
  \label{eq:flight-2d-FM}
\end{equation}
Neither force has a $y$ component, so $v_y$ stays zero and the ball stays in the plane.
Define the drag constant $\kappa = \rho A/(2m)$ (\unit{1/m}) and the backspin sign
$s_b = -\operatorname{sgn}(\omega_y)$ ($s_b = +1$ for backspin toward $+x$).
Dividing Eq.~\eqref{eq:flight-eom-v} by $m$ and inserting Eqs.~\eqref{eq:flight-drag} and \eqref{eq:flight-2d-FM}:
\begin{align}
  \frac{dv_x}{dt} &= -\kappa\,C_d\,|\vect{v}_{\mathrm{rel}}|\,v_x - s_b\,\kappa\,C_L(S)\,|\vect{v}_{\mathrm{rel}}|\,v_z, \label{eq:flight-2d-ax}\\
  \frac{dv_z}{dt} &= -g - \kappa\,C_d\,|\vect{v}_{\mathrm{rel}}|\,v_z + s_b\,\kappa\,C_L(S)\,|\vect{v}_{\mathrm{rel}}|\,v_x, \label{eq:flight-2d-az}\\
  \omega_y(t) &= \omega_{y,0}\,e^{-kt}, \qquad S = \frac{|\omega_y|\,R}{|\vect{v}_{\mathrm{rel}}|}. \label{eq:flight-2d-w}
\end{align}
These are the equations in \code{feasibility.integrate\_batch} (with $s_b$ given by the sign of the
signed exit spin parameter $S_0$, positive for backspin) and in \code{contact.\_acc}
(written with $\omega_y$ directly, backspin negative).
The nominal value is $\kappa = 0.0952$\,\unit{1/m}\derived.
% Code note: feasibility.integrate_batch has no wind (v_rel = v) and leaves eps_v out of the Magnus
% direction denominator (it keeps eps_v only inside S). contact._acc includes eps_v in both. The
% difference is below 1e-7 m on the worked shot. Sign conventions differ between modules: feasibility
% uses a signed spin magnitude (> 0 = backspin), contact and flight use omega_y (< 0 = backspin toward +x).

\paragraph{Worked example.}
At launch ($v_x = 3.5$\assumed, $v_z = 6.062$\,\unit{m/s}\derived, $|\vect{v}| = 7$\,\unit{m/s}\assumed,
$C_L = 0.25$\derived), the Magnus terms of Eqs.~\eqref{eq:flight-2d-ax}--\eqref{eq:flight-2d-az} are
$-0.0952\times0.25\times7\times6.062 = -1.010$\,\unit{m/s\textsuperscript{2}}\derived\ in $x$ and
$+0.0952\times0.25\times7\times3.5 = +0.583$\,\unit{m/s\textsuperscript{2}}\derived\ in $z$.
Multiplying by $m$ gives $(-0.0252, 0.0146)$\,\unit{N}\derived, the same as the 3D \code{aero\_forces}
result in Section~\ref{sec:flight-magnus}.

\paragraph{Assumptions.}
\begin{itemize}
  \item Pure backspin about a horizontal axis perpendicular to the launch plane.
  \item No sideways wind and no sideways launch component.
\end{itemize}

% ---------------------------------------------------------------------------------------------
\subsection{How large are drag and Magnus compared with gravity?}
\label{sec:flight-scale}

\paragraph{Motivation.}
Before trusting any trajectory map it helps to know the size of each force.
Ratios to $g$ make that easy to read.

\paragraph{Terminal speed.}
For a ball falling straight down without spin, the speed stops growing when drag balances weight.
Setting $\tfrac12\rho C_d A v_t^2 = mg$ and solving for the terminal speed $v_t$ (\unit{m/s}):
\begin{equation}
  v_t = \sqrt{\frac{2\,m\,g}{\rho\,C_d\,A}}.
  \label{eq:flight-vt}
\end{equation}
For $C_d = 0.33$\assumed, $0.5$\assumed\ and $0.9$\assumed\ this gives
$v_t = 17.67$\derived, $14.35$\derived\ and $10.70$\,\unit{m/s}\derived.
% Code vs docs: docs/02_aero.md section 3 quotes 14.6 m/s, computed from the older aero_params.json
% nominal Cd = 0.48. With params_lit.json Cd = 0.5 it is 14.35 m/s.

\paragraph{Drag and Magnus over $g$.}
Divide each force magnitude by $mg$ and use Eq.~\eqref{eq:flight-vt} to replace $\rho A/(2m)$ by
$g/(C_d v_t^2)$:
\begin{align}
  \frac{a_D}{g} &= \frac{\rho\,C_d\,A\,|\vect{v}_{\mathrm{rel}}|^2}{2\,m\,g}
               = \left(\frac{|\vect{v}_{\mathrm{rel}}|}{v_t}\right)^{2}, \label{eq:flight-aD}\\
  \frac{a_M}{g} &= \frac{\rho\,C_L\,A\,|\vect{v}_{\mathrm{rel}}|^2}{2\,m\,g}
               = \frac{C_L}{C_d}\left(\frac{|\vect{v}_{\mathrm{rel}}|}{v_t}\right)^{2}, \label{eq:flight-aM}\\
  \frac{a_M}{a_D} &= \frac{C_L}{C_d}. \label{eq:flight-ratio}
\end{align}
Here $a_D$ and $a_M$ (\unit{m/s\textsuperscript{2}}) are the drag and Magnus accelerations.
Both grow with the square of the airspeed, so both are small near the apex and largest at launch and
at the end of a long fall.
Because $C_L \le C_{L,\max} \le 0.4$\assumed\ and $C_d \ge 0.33$\assumed, Magnus is smaller than drag
for every set except the high-lift, low-drag corner ($0.4/0.33 = 1.2$\derived).
At the nominal values the ratio is at most $0.25/0.5 = 0.5$\derived.

\begin{table}[htbp]
  \centering
  \caption{Accelerations as a fraction of $g$ (nominal $\rho$, $m$, $R$). Drag for the three $C_d$
    bounds; Magnus at $S = 1$\assumed\ for the four lift sets (only the floor set is below its cap there).}
  \label{tab:flight-scale}
  \begin{tabular}{@{}l lll llll@{}}
    \toprule
    & \multicolumn{3}{c}{$a_D/g$} & \multicolumn{4}{c}{$a_M/g$ at $S=1$} \\
    \cmidrule(lr){2-4}\cmidrule(l){5-8}
    $|\vect{v}_{\mathrm{rel}}|$ (\unit{m/s}) & $C_d$ 0.33 & 0.5 & 0.9 & floor & low & nominal & high \\
    \midrule
    3\assumed & 0.029\derived & 0.044\derived & 0.079\derived & 0.002\derived & 0.013\derived & 0.022\derived & 0.035\derived \\
    5\assumed & 0.080\derived & 0.121\derived & 0.218\derived & 0.005\derived & 0.036\derived & 0.061\derived & 0.097\derived \\
    7\assumed & 0.157\derived & 0.238\derived & 0.428\derived & 0.010\derived & 0.071\derived & 0.119\derived & 0.190\derived \\
    9\assumed & 0.260\derived & 0.393\derived & 0.708\derived & 0.016\derived & 0.118\derived & 0.197\derived & 0.315\derived \\
    \bottomrule
  \end{tabular}
\end{table}

\paragraph{Path bending.}
What changes the shape of the path is the force component perpendicular to the velocity.
For a path with local radius of curvature $r_c$ (\unit{m}), the normal acceleration is
\begin{equation}
  a_\perp = \frac{|\vect{v}|^2}{r_c}
  \quad\Longrightarrow\quad
  r_c = \frac{|\vect{v}|^2}{a_\perp}.
  \label{eq:flight-curv}
\end{equation}
Drag has no normal component.
Gravity contributes $-g\cos\theta$, and in still air the Magnus force is entirely normal
(Eq.~\eqref{eq:flight-FM-perp}).

\paragraph{Worked example.}
At launch ($7$\,\unit{m/s}\assumed, $60^\circ$\assumed, nominal, $S_0 = 1.06$\assumed):
$a_D/g = 0.238$\derived\ and $a_M/g = 0.119$\derived.
The normal accelerations are $-4.903$\,\unit{m/s\textsuperscript{2}}\derived\ from gravity and
$+1.166$\,\unit{m/s\textsuperscript{2}}\derived\ from Magnus, so backspin cancels $24$\,\%\derived\ of the downward
bending at that moment.
The radius of curvature grows from $9.99$\derived\ to $13.11$\,\unit{m}\derived.
At the apex of the same shot ($t = 0.594$\,\unit{s}\predicted), $v_z = 0$ but the ball is still moving
horizontally at $v_x = 2.885$\,\unit{m/s}\predicted; it is not stationary.
There $S = 2.53$\predicted\ but $C_L$ is still capped at $0.25$\derived, and the Magnus acceleration is only
$0.020\,g$\predicted, pointing straight up.
Spin has its largest effect on the fast parts of the path, not at the apex.

\paragraph{Wind example.}
A $1$\,\unit{m/s}\assumed\ headwind, $\vect{u} = (-1, 0, 0)$\,\unit{m/s}\assumed, makes
$\vect{v}_{\mathrm{rel}} = (4.5, 0, 6.062)$\,\unit{m/s}\derived\ with $|\vect{v}_{\mathrm{rel}}| = 7.55$\,\unit{m/s}\derived.
The spin parameter falls to $S = 0.983$\derived\ (still above $S_{\mathrm{sat}}$), and
\code{aero\_forces} gives $|\vect{F}_D| = 0.0677$\,\unit{N}\predicted\ and
$|\vect{F}_M| = 0.0339$\,\unit{N}\predicted, both $16$\,\%\derived\ larger than in still air because
$|\vect{v}_{\mathrm{rel}}|^2$ grew from $49$\derived\ to $57$\,\unit{m\textsuperscript{2}/s\textsuperscript{2}}\derived.

\paragraph{Assumptions.}
\begin{itemize}
  \item The table uses the nominal $\rho$, $m$ and $R$; the bounds of Table~\ref{tab:flight-params}
    scale every entry by $\rho R^2/m$, which moves it by $-27$\,\%\derived\ to $+14$\,\%\derived.
  \item The terminal speed is for a non-spinning ball falling vertically; the shots in this project
    stay well below it.
\end{itemize}

% ---------------------------------------------------------------------------------------------
\subsection{Numerical integration}
\label{sec:flight-numerics}

\paragraph{Motivation.}
Equations~\eqref{eq:flight-eom-r}--\eqref{eq:flight-eom-w} have no closed-form solution once drag and lift
are both present, so they are integrated numerically.
The project uses two kinds of integrator, and we need to know that they agree.

\paragraph{Adaptive Runge--Kutta in \code{flight.simulate}.}
\code{simulate} calls SciPy's \code{solve\_ivp} with the Dormand--Prince 5(4) method (\code{RK45};
\code{DOP853} is an option).
At every step the method estimates its own local error $e_i$ for each state component $y_i$ and
accepts the step only if
\begin{equation}
  \sqrt{\frac{1}{9}\sum_{i=1}^{9}\left(\frac{e_i}{a_{\mathrm{tol}} + r_{\mathrm{tol}}\,\max(|y_{n,i}|,|y_{n+1,i}|)}\right)^{2}} \le 1 ,
  \label{eq:flight-tol}
\end{equation}
otherwise it retries with a smaller step.
The defaults are $r_{\mathrm{tol}} = 10^{-9}$\assumed, $a_{\mathrm{tol}} = 10^{-11}$\assumed\ and a maximum
step of $10^{-3}$\,\unit{s}\assumed.
Dense output gives the state at any time.
Event functions find the apex (first crossing of $v_z$ from positive to negative, non-terminal),
the floor ($z$ falling through \code{z\_ground}, terminal) and optional planes.
The apex event records $v_z = 0$ together with the nonzero $v_x$.

\paragraph{Fixed-step RK4 in the sweeps.}
The sweeps integrate many shots at once with the classical fourth-order Runge--Kutta method.
For $\dot{\vect{y}} = \vect{f}(t,\vect{y})$ and step $h$ (\unit{s}):
\begin{align}
  \vect{k}_1 &= \vect{f}(t_n,\vect{y}_n), \qquad
  \vect{k}_2 = \vect{f}\bigl(t_n+\tfrac h2,\vect{y}_n+\tfrac h2\vect{k}_1\bigr), \nonumber\\
  \vect{k}_3 &= \vect{f}\bigl(t_n+\tfrac h2,\vect{y}_n+\tfrac h2\vect{k}_2\bigr), \qquad
  \vect{k}_4 = \vect{f}(t_n+h,\vect{y}_n+h\vect{k}_3), \nonumber\\
  \vect{y}_{n+1} &= \vect{y}_n + \tfrac{h}{6}\bigl(\vect{k}_1 + 2\vect{k}_2 + 2\vect{k}_3 + \vect{k}_4\bigr). \label{eq:flight-rk4}
\end{align}
\code{feasibility.integrate\_batch} and \code{robustness.integrate3d} write the position update in a
shorter form.
Since $d\vect{r}/dt = \vect{v}$, the four position slopes are $\vect{v}_n$,
$\vect{v}_n+\tfrac h2\vect{a}_1$, $\vect{v}_n+\tfrac h2\vect{a}_2$ and $\vect{v}_n+h\vect{a}_3$, where
$\vect{a}_j$ are the acceleration stages.
Inserting them into Eq.~\eqref{eq:flight-rk4}:
\begin{equation}
  \vect{r}_{n+1}
  = \vect{r}_n + \tfrac h6\bigl(6\vect{v}_n + h\vect{a}_1 + h\vect{a}_2 + h\vect{a}_3\bigr)
  = \vect{r}_n + h\Bigl(\vect{v}_n + \tfrac h6(\vect{a}_1+\vect{a}_2+\vect{a}_3)\Bigr),
  \label{eq:flight-rk4-pos}
\end{equation}
which is exactly the code's \code{x + dt*(vx + dt/6*(a1x + a2x + a3x))}.
The global error of RK4 scales as
\begin{equation}
  \bigl|\vect{r}_h(t) - \vect{r}(t)\bigr| \approx C\,h^{4},
  \label{eq:flight-rk4-order}
\end{equation}
so halving $h$ should cut the error by $2^4 = 16$.
The sweeps use $h = 2$\,\unit{ms}\assumed\ (\code{integrate\_batch}, \code{integrate3d}) and
$h = 1$\,\unit{ms}\assumed\ (\code{contact.free\_flight}).
The two batch integrators store positions in single precision (\code{float32}).

\paragraph{Worked check: do they agree?}
Take the worked shot ($7$\,\unit{m/s}\assumed, $60^\circ$\assumed, $S_0 = 1.06$\assumed, nominal set,
$1.2$\,\unit{s}\assumed\ of flight) and use \code{DOP853} with $r_{\mathrm{tol}} = 10^{-12}$\assumed\ as the
reference.
\begin{itemize}
  \item \code{simulate} (RK45, default tolerances): largest position difference
    $1.1\times10^{-15}$\,\unit{m}\predicted.
  \item \code{contact.free\_flight} (RK4, double precision), error at $t = 1.2$\,\unit{s}:
    $9.3\times10^{-11}$\predicted, $5.8\times10^{-12}$\predicted, $3.5\times10^{-13}$\predicted\ and
    $2.2\times10^{-14}$\,\unit{m}\predicted\ for $h = 8$\assumed, $4$\assumed, $2$\assumed\ and
    $1$\,\unit{ms}\assumed.
    Successive ratios are $16.1$\derived, $16.6$\derived\ and $16.0$\derived, as
    Eq.~\eqref{eq:flight-rk4-order} predicts.
  \item \code{integrate\_batch} (2D) and \code{integrate3d} (3D), largest position difference:
    $1.2\times10^{-7}$\,\unit{m}\predicted\ for every step from $8$\assumed\ to $1$\,\unit{ms}\assumed.
    This floor is the \code{float32} rounding of positions between $2$\derived\ and $4$\,\unit{m}\derived\
    (spacing $2.4\times10^{-7}$\,\unit{m}\derived, so a rounding error of at most
    $1.2\times10^{-7}$\,\unit{m}\derived), not truncation error.
    % Fixed: the draft gave 1.3e-7 m and a spacing of 4.8e-7 m "near 4 m"; this shot (launched from the
    % origin, as integrate_batch requires) reaches at most 3.54 m in 1.2 s, where the float32 spacing is
    % 2.4e-7 m. Recomputed 2026-10-30: 1.17e-7 to 1.19e-7 m for h = 8 to 1 ms.
\end{itemize}
The unit tests demand agreement better than $10^{-5}$\,\unit{m}\assumed\ (\code{tests/test\_hive.py})
and $2\times10^{-6}$\,\unit{m}\assumed\ (\code{tests/test\_contact.py}).
Integration error is therefore at least four orders of magnitude smaller than the millimetre- to
centimetre-scale differences between lift sets in Table~\ref{tab:flight-setfx}.
The parameters, not the numerics, limit what the model can say.

\paragraph{Assumptions.}
\begin{itemize}
  \item The check is for one smooth free-flight path; contact events are handled separately
    (Section on contact) and have their own step control.
  \item Agreement between integrators shows the equations are solved correctly; it says nothing about
    whether the equations describe a real POLLEN ball.
\end{itemize}

% ---------------------------------------------------------------------------------------------
\subsection{Summary of modelling assumptions}
\label{sec:flight-assumptions}

\begin{enumerate}
  \item Rigid sphere, point-mass translation, inertial frame, flat floor; buoyancy and added mass
    neglected.
  \item Forces use the air-relative velocity $\vect{v}-\vect{u}$; the wind is prescribed.
  \item Drag: $\vect{F}_D = -\tfrac12\rho C_d A|\vect{v}_{\mathrm{rel}}|\vect{v}_{\mathrm{rel}}$
    with $C_d$ constant per run (bounds $0.33$\assumed\ to $0.9$\assumed); no spin-dependent drag.
  \item Magnus: direction $\operatorname{unit}(\vect{\omega}\times\vect{v}_{\mathrm{rel}})$, magnitude
    $\tfrac12\rho C_L A|\vect{v}_{\mathrm{rel}}|^2$; no reverse Magnus, no non-Magnus hole forces.
  \item Lift law $C_L = \min(c_1 S, C_{L,\max})$, independent of $\mathrm{Re}$. Above
    $S_{\mathrm{sat}} = C_{L,\max}/c_1$ extra spin changes nothing; this is a model assumption.
  \item Spin decay $\vect{\omega} = \vect{\omega}_0 e^{-kt}$ with constant $k$; the spin axis is fixed.
  \item Backspin axis $\operatorname{unit}(\vect{v}_0\times\hat{\vect{z}})$, so pure backspin keeps the
    ball in the vertical launch plane.
  \item All aerodynamic coefficients come from assumed, literature-based ranges (pickleball and
    smooth-sphere data). They are not proven bounds and not POLLEN measurements. The worked numbers
    above are single model runs and say nothing about shot-to-shot repeatability.
\end{enumerate}
